% (c) Nikita Lisitsa, lisyarus@gmail.com, 2026

\documentclass[10pt]{beamer}

\usepackage[T2A]{fontenc}
\usepackage[russian]{babel}
\usepackage{minted}

\usepackage[most]{tcolorbox}
\tcbuselibrary{minted}

\usepackage{graphicx}
\graphicspath{ {./images/} }

\usepackage{adjustbox}

\usepackage{color}
\usepackage{soul}

\usepackage{hyperref}

\usetheme{metropolis}

\definecolor{red}{rgb}{1,0,0}
\definecolor{codebg}{RGB}{29,35,49}
\setminted{fontsize=\footnotesize}

\makeatletter
\newcommand{\slideimage}[1]{
  \begin{figure}
    \begin{adjustbox}{width=\textwidth, totalheight=\textheight-2\baselineskip-2\baselineskip,keepaspectratio}
      \includegraphics{#1}
    \end{adjustbox}
  \end{figure}
}
\makeatother

\title{Язык программирования C++}
\subtitle{Лекция 11: временные объекты, new/delete с объектами, умные указатели, наследование}
\date{2026}
\author{Никита Лисица}

\setbeamertemplate{footline}[frame number]

\begin{document}

\frame{\titlepage}

\begin{frame}[fragile]
\frametitle{Временные объекты}
\begin{itemize}
\item Из-за присутствия конструкторов и деструкторов, в C++ приходится внимательнее следить за временем жизни объектов
\pause
\usemintedstyle{lightbulb}
\begin{tcblisting}{listing engine=minted, minted language=cpp, minted options={fontsize=\scriptsize}, listing only, colback=codebg, colframe=codebg, rounded corners}
int max(int a, int b) {
  return a < b ? b : a;
}

int main() {
  // `a` живёт от этой строки и до конца main
  int a = 2;

  // `b` живёт от этой строки и до конца main
  int b = 3;

  max(a, b);
}
\end{tcblisting}
\end{itemize}
\end{frame}

\begin{frame}[fragile]
\frametitle{Временные объекты}
\begin{itemize}
\usemintedstyle{tango}
\item Что происходит, если написать \mintinline{cpp}|max(2, 3)|? \pause Создаются временные объекты:
\pause
\usemintedstyle{lightbulb}
\begin{tcblisting}{listing engine=minted, minted language=cpp, minted options={fontsize=\scriptsize}, listing only, colback=codebg, colframe=codebg, rounded corners}
int max(int a, int b) {
  return a < b ? b : a;
}

int main() {
  // Здесь, `2` и `3` это временные объекты
  // Они создаются перед вызовом функции, и
  // живут до окончания выполнения всего выражения
  // Возвращаемое значение - тоже временный объект
  max(2, 3);
}
\end{tcblisting}
\end{itemize}
\end{frame}

\begin{frame}[fragile]
\frametitle{Временные объекты}
\begin{itemize}
\usemintedstyle{tango}
\item Это не очень важно для встроенных типов, но важно для классов с нетривиальными конструкторами:
\pause
\usemintedstyle{lightbulb}
\begin{tcblisting}{listing engine=minted, minted language=cpp, minted options={fontsize=\scriptsize}, listing only, colback=codebg, colframe=codebg, rounded corners}
BigInt max(const BigInt& a, const BigInt& b) {
  return a < b ? b : a;
}

int main() {
  // В этой строке создаются два временных объекта
  // типа BigInt
  max(
    BigInt("12347987955"),
    BigInt("867577959256")
  );
}
\end{tcblisting}
\end{itemize}
\end{frame}

\begin{frame}[fragile]
\frametitle{Параметры по умолчанию}
\begin{itemize}
\usemintedstyle{tango}
\item У функций можно указывать параметры по умолчанию, тогда их не обязательно передавать:
\usemintedstyle{lightbulb}
\begin{tcblisting}{listing engine=minted, minted language=cpp, minted options={fontsize=\scriptsize}, listing only, colback=codebg, colframe=codebg, rounded corners}
void startServer(const char* name, int port = 8080,
  const char* host = "0.0.0.0");

int main()
{
  startServer("minecraft", 25565, "127.0.0.1");
  startServer("minecraft", 5432);
  startServer("ping");
}
\end{tcblisting}
\end{itemize}
\end{frame}

\begin{frame}[fragile]
\frametitle{Параметры по умолчанию}
\begin{itemize}
\usemintedstyle{tango}
\item Они могут быть только последними параметрами (иначе непонятно, какие параметры передали)
\pause
\item Вместе с перегрузкой функций может приводить к неожиданным эффектам или к неопределённости
\pause
\item Параметры по умолчанию подставляются компилятором в код вызова функции \(\Longrightarrow\) должны быть описаны в объявлении функции (обычно в header-файле)
\end{itemize}
\end{frame}

\begin{frame}[fragile]
\frametitle{new/delete с объектами}
\begin{itemize}
\usemintedstyle{tango}
\item Пусть, у нас есть класс
\usemintedstyle{lightbulb}
\begin{tcblisting}{listing engine=minted, minted language=cpp, minted options={fontsize=\scriptsize}, listing only, colback=codebg, colframe=codebg, rounded corners}
class Server {
public:
  Server(int port);
  ~Server();
};
\end{tcblisting}
\pause
\item Как нам создать объект такого класса в куче?
\pause
\usemintedstyle{tango}
\item \mintinline{cpp}|malloc| позволит выделить память, но не вызовет конструктор
\pause
\item \mintinline{cpp}|free| не вызовет деструктор
\end{itemize}
\end{frame}

\begin{frame}[fragile]
\frametitle{new/delete с объектами}
\begin{itemize}
\usemintedstyle{tango}
\item Решение: \mintinline{cpp}|new/delete|!
\usemintedstyle{lightbulb}
\begin{tcblisting}{listing engine=minted, minted language=cpp, minted options={fontsize=\scriptsize}, listing only, colback=codebg, colframe=codebg, rounded corners}
// Можем передать параметры конструктора
Server * server = new Server(5432);
...
// Вызывает деструктор
delete server;
\end{tcblisting}
\pause
\usemintedstyle{tango}
\item Почему компилятор сам не вставит вызов деструктора? \pause Он не знает, какое время жизни мы хотим у объекта \mintinline{cpp}|*server|
\pause
\item Он вызовет деструктор у самого указателя \mintinline{cpp}|server| -- но деструкторы примитивных типов ничего не делают
\end{itemize}
\end{frame}

\begin{frame}[fragile]
\frametitle{Placement new}
\begin{itemize}
\usemintedstyle{tango}
\item Можно и через \mintinline{cpp}|malloc/free|, но руками вызвать конструктор и деструктор:
\usemintedstyle{lightbulb}
\begin{tcblisting}{listing engine=minted, minted language=cpp, minted options={fontsize=\scriptsize}, listing only, colback=codebg, colframe=codebg, rounded corners}
Server * server = malloc(sizeof(Server));

// Особый синтаксис вызова конструктора
// у объекта по указателю (не выделяет память)
// - т.н. placement new
new (server) Server(5432);
...

// Руками вызываем деструктор
server->~Server();

free(server);
\end{tcblisting}
\end{itemize}
\end{frame}

\begin{frame}[fragile]
\frametitle{Placement new}
\begin{itemize}
\usemintedstyle{tango}
\item Обычно нужно только при реализации структур данных (напр. \mintinline{cpp}|std::vector|), в большинстве случаев хватает обычных \mintinline{cpp}|new/delete|
\pause
\item Даже если вы выделили объект через \mintinline{cpp}|malloc| + placement new, его нельзя удалять через \mintinline{cpp}|delete|, только через ручной вызов деструктора + \mintinline{cpp}|free|
\end{itemize}
\end{frame}

\begin{frame}[fragile]
\frametitle{new[]/delete[] с объектами}
\begin{itemize}
\usemintedstyle{tango}
\item Аналогично, \mintinline{cpp}|new[]| не только выделит память, но и вызовет конструктор по умолчанию у всех объектов
\pause
\item \mintinline{cpp}|delete[]| пройдётся по всем объектам и вызовет у каждого деструктор
\pause
\usemintedstyle{lightbulb}
\begin{tcblisting}{listing engine=minted, minted language=cpp, minted options={fontsize=\scriptsize}, listing only, colback=codebg, colframe=codebg, rounded corners}
// Вызовет конструктор BigInt() у элементов
// от bigInts[0] до bigInts[99]
BigInt * bigInts = new BigInt[100];

// Вызовет деструктор у элементов
// от bigInts[0] до bigInts[99]
delete[] bigInts;
\end{tcblisting}
\pause
\usemintedstyle{tango}
\item А если нет конструктора по умолчанию? \pause \mintinline{cpp}|malloc| + placement new
\end{itemize}
\end{frame}

\begin{frame}[fragile]
\frametitle{Умные указатели}
\begin{itemize}
\usemintedstyle{tango}
\item Многовато вызовов \mintinline{cpp}|delete|! Воспользуемся инкапсуляцией и деструкторами \(\Longrightarrow\) \textit{умные указатели}
\pause
\item Без умных указателей:
\usemintedstyle{lightbulb}
\begin{tcblisting}{listing engine=minted, minted language=cpp, minted options={fontsize=\scriptsize}, listing only, colback=codebg, colframe=codebg, rounded corners}
Object* p = new Object(...);
p->doSomething();
// Надо не забыть вызвать delete
delete p;
\end{tcblisting}
\pause
\item С умными указателями:
\begin{tcblisting}{listing engine=minted, minted language=cpp, minted options={fontsize=\scriptsize}, listing only, colback=codebg, colframe=codebg, rounded corners}
scoped_ptr p = new Object(...);
p->doSomething();
// деструктор scoped_ptr вызовет delete
\end{tcblisting}
\end{itemize}
\end{frame}

\begin{frame}[fragile]
\frametitle{scoped\_ptr}
\usemintedstyle{lightbulb}
\begin{tcblisting}{listing engine=minted, minted language=cpp, minted options={fontsize=\scriptsize}, listing only, colback=codebg, colframe=codebg, rounded corners}
class scoped_ptr {
private:
  Object* ptr_;

public:
  scoped_ptr(Object* ptr);
  ~scoped_ptr();

  Object* get() const;
  Object& operator * () const;
  Object* operator -> () const;
  explicit operator bool() const;

  // Запрет копирования после C++11:
  scoped_ptr(const scoped_ptr& other) = delete;
  const scoped_ptr& operator = (const scoped_ptr& p) = delete;

  // Запрет копирования до C++11:
private:
  scoped_ptr(const scoped_ptr& other);
  const scoped_ptr& operator = (const scoped_ptr& p);
};
\end{tcblisting}
\end{frame}

\begin{frame}[fragile]
\frametitle{scoped\_ptr}
\usemintedstyle{lightbulb}
\begin{tcblisting}{listing engine=minted, minted language=cpp, minted options={fontsize=\scriptsize}, listing only, colback=codebg, colframe=codebg, rounded corners}
scoped_ptr::scoped_ptr(Object* ptr)
  : ptr_(ptr)
{}

scoped_ptr::~scoped_ptr() {
  delete ptr_;
}

Object* scoped_ptr::get() const { return ptr_; }

Object& scoped_ptr::operator * () const { return *ptr_; }

Object* scoped_ptr::operator -> () const { return ptr_; }

scoped_ptr::operator bool() const { return ptr_ != nullptr; }
\end{tcblisting}
\end{frame}

\begin{frame}[fragile]
\frametitle{scoped\_ptr}
\begin{itemize}
\usemintedstyle{tango}
\item Зачем запрет копирования? \pause Чтобы два разных \mintinline{cpp}|scoped_ptr| не пытались удалить один и тот же указатель
\pause
\item Как работает \mintinline{cpp}|operator ->|? \pause Возвращает либо указатель (у которого есть встроенный оператор), либо что-нибудь, у чего тоже есть \mintinline{cpp}|operator ->|, и вызывает его рекурсивно
\end{itemize}
\end{frame}

\begin{frame}[fragile]
\frametitle{Умные указатели}
\begin{itemize}
\usemintedstyle{tango}
\item \mintinline{cpp}|scoped_ptr|: нельзя копировать, нельзя передать владение объектом (напр. в функцию или из функции), объект живёт ровно внутри одного scope, только один \mintinline{cpp}|scoped_ptr| указывает на объект
\pause
\item \mintinline{cpp}|unique_ptr|: нельзя копировать, можно перемещать (\textit{move}, об этом позже), можно передавать владение, только один \mintinline{cpp}|unique_ptr| указывает на объект
\pause
\item \mintinline{cpp}|shared_ptr|: можно копировать (\(\Longrightarrow\) можно передавать владение), несколько \mintinline{cpp}|shared_ptr| могут указывать на один объект
\pause
\item \textbf{NB:} На самом деле, умные указатели -- это шаблоны классов, но о шаблонах мы поговорим позже
\end{itemize}
\end{frame}

\begin{frame}[fragile]
\frametitle{shared\_ptr}
\begin{itemize}
\usemintedstyle{tango}
\item Идея: храним счётчик ссылок
\pause
\item Когда \mintinline{cpp}|shared_ptr| копируется, увеличиваем счётчик
\pause
\item Когда \mintinline{cpp}|shared_ptr| удаляется, уменьшаем счётчик
\pause
\item Когда счётчик доходит до нуля (т.е. удалили последний \mintinline{cpp}|shared_ptr|, ссылавшийся на этот объект), удаляем сам объект
\pause
\item \textbf{NB:} Не подходит, если есть циклические ссылки!
\end{itemize}
\end{frame}

\begin{frame}[fragile]
\frametitle{shared\_ptr}
\begin{itemize}
\usemintedstyle{tango}
\item Счётчик ссылок должен быть общим между всеми копиями \mintinline{cpp}|shared_ptr|, указывающими на один объект \(\Longrightarrow\) его придётся тоже хранить в куче
\usemintedstyle{lightbulb}
\begin{tcblisting}{listing engine=minted, minted language=cpp, minted options={fontsize=\scriptsize}, listing only, colback=codebg, colframe=codebg, rounded corners}
class shared_ptr {
private:
  struct storage {
    std::size_t refcount = 0;
    Object* pointer = nullptr;
  };

  storage* storage_;
  
public:
  ...
};
\end{tcblisting}
\end{itemize}
\end{frame}

\begin{frame}[fragile]
\frametitle{shared\_ptr}
\begin{itemize}
\usemintedstyle{tango}
\item Настоящий \mintinline{cpp}|shared_ptr| устроен сложнее
\pause
\item Позволяет передать произвольную функцию удаления (по умолчанию делает просто \mintinline{cpp}|delete|)
\pause
\item Позволяет ссылаться на часть объекта
\pause
\item Позволяет разбивать циклические ссылки или узнавать об удалении объекта через \mintinline{cpp}|weak_ptr|
\end{itemize}
\end{frame}

\begin{frame}[fragile]
\frametitle{Наследование}
\begin{itemize}
\usemintedstyle{tango}
\item Одна из основных концепций ООП
\pause
\item Хотим переиспользовать код классов в других классах
\pause
\item Вариант №1 -- \textit{композиция}: один класс является полем другого
\pause
\item Вариант №2 -- \textit{наследование}: один класс строится на базе другого
\end{itemize}
\end{frame}

\begin{frame}[fragile]
\frametitle{Наследование}
\begin{itemize}
\usemintedstyle{tango}
\usemintedstyle{lightbulb}
\begin{tcblisting}{listing engine=minted, minted language=cpp, minted options={fontsize=\scriptsize}, listing only, colback=codebg, colframe=codebg, rounded corners}
class Button {
private:
  const char* caption_;
  Action clickAction_;

public:
  Button(const char* caption)
    : caption_(caption)
  {}

  void draw() {
    drawRectangle(Color::Gray);
    drawText(caption_);
  }

  void onClick(Action clickAction) {
    clickAction_ = clickAction;
  }

  void click() {
    clickAction_.execute();
  }
};
\end{tcblisting}
\end{itemize}
\end{frame}

\begin{frame}[fragile]
\frametitle{Наследование}
\begin{itemize}
\usemintedstyle{tango}
\usemintedstyle{lightbulb}
\begin{tcblisting}{listing engine=minted, minted language=cpp, minted options={fontsize=\scriptsize}, listing only, colback=codebg, colframe=codebg, rounded corners}
class ColoredButton : public Button {
private:
  Color color_;

public:
  ColoredButton(const char* caption, Color color)
    : Button(caption) // вызов конструктора
    , color_(color)
  {}

  void draw() {
    drawRectangle(color);
    // Ошибка: caption_ - private поле, об этом позже
    drawText(caption_);
  }

  void click() {
    color_ = Color::Red;
    // Вызов метода из базового класса
    Button::click();
  }
};
\end{tcblisting}
\end{itemize}
\end{frame}

\begin{frame}[fragile]
\frametitle{Наследование}
\begin{itemize}
\usemintedstyle{tango}
\item \mintinline{cpp}|Button| -- \textit{базовый класс (base class)}
\pause
\item \mintinline{cpp}|ColoredButton| -- \textit{производный класс (derived class)}, или \textit{класс-наследник}
\pause
\item В \mintinline{cpp}|ColoredButton| добавляются все поля и методы из базового класса
\pause
\item Конструктор класса-наследника должен вызвать конструктор базового класса (как и с полями, по умолчанию вызовется конструктор без аргументов)
\pause
\item Деструктор класса-наследника \underline{автоматически} вызовет деструктор базового класса после того, как вызовет деструкторы всех полей класса-наследника
\pause
\item Указатель на класс-наследник автоматически приводится к типу указателя на базовый класс
\end{itemize}
\end{frame}

\begin{frame}[fragile]
\frametitle{Наследование}
\begin{itemize}
\usemintedstyle{tango}
\item Хотим обрабатывать разные объекты одним и тем же кодом:
\usemintedstyle{lightbulb}
\begin{tcblisting}{listing engine=minted, minted language=cpp, minted options={fontsize=\scriptsize}, listing only, colback=codebg, colframe=codebg, rounded corners}
Button * buttons[3];
buttons[0] = new Button("Save");
buttons[1] = new Button("Load");
buttons[2] = new ColoredButton("Quit", Color::Blue);

void drawUI() {
  for (int i = 0; i < 3; ++i)
    buttons[i]->draw();
}
\end{tcblisting}
\pause
\usemintedstyle{tango}
\item Функция \mintinline{cpp}|drawUI| не знает, какие классы на самом деле лежат по разным указателям, и всегда будет вызывать \mintinline{cpp}|Button::draw()| из базового класса
\pause
\item Решение: \textit{виртуальные методы}
\end{itemize}
\end{frame}

\begin{frame}[fragile]
\frametitle{Виртуальные методы}
\begin{itemize}
\usemintedstyle{lightbulb}
\begin{tcblisting}{listing engine=minted, minted language=cpp, minted options={fontsize=\scriptsize}, listing only, colback=codebg, colframe=codebg, rounded corners}
class Button {
public:
  virtual void draw();
  virtual void click();
  ...
};
\end{tcblisting}
\pause
\usemintedstyle{tango}
\item \mintinline{cpp}|virtual| говорит компилятору, что при вызове этого метода надо посмотреть, какой на самом деле класс лежит по этому указателю, и вызвать правильный метод
\pause
\item Теперь \mintinline{cpp}|buttons[2]->draw()| вызовет \mintinline{cpp}|ColoredButton::draw()|!
\pause
\item Как это работает, обсудим в следующий раз
\end{itemize}
\end{frame}

\end{document}